\documentclass{article}
\usepackage{graphicx}
\usepackage{wrapfig}
%\usepackage{inconsolata}
\usepackage{enumerate}
\usepackage{hyperref}
\usepackage{verbatim}
\usepackage[parfill]{parskip}
\usepackage[margin = 2.5cm]{geometry}

\usepackage[T1]{fontenc}


\begin{document}

\title{AWS\\ IN720 Virtualisation}
\date{}
\maketitle

\section*{Introduction}
\emph{Amazon Web Services} (AWS) are probably the leading (although not the only) public cloud services today. In this lab we will use AWS' EC2 service to launch a simple web site.

\section{Launch an EC2 instance}

Carry out the following procedure to launch your Amazon EC2 instance:

\begin{enumerate}

\item Sign into the AWS web console at \url{https://355631893594.signin.aws.amazon.com/console} using credentials given to you by the lecturer.

\item From the services menu, select EC2.

\item Set your region to US West (Northern California) at the top right of the page. Then click ``Launch Instance''.

\item From the image selection menu, choose Ubuntu 14.04.

\item From the instance type menu, choose t2.micro.  Click next.

\item At the instance details page, keep the defaults. Click next.

\item From the storage screen, keep the default root volume.

\item At the tagging screen, set a tag with key ``Name''. Use your surname as the value.  Click next.

\item At the security group page, create a new group. Create firewall rules that allow SSH and HTTP access from anywhere. Click Review and Launch.

\item Review the settings and click launch.  You'll be prompted to choose a key pair. Create a new key at this stage and download the key file.

\item As your instance launches, find its entry in the instance list and note its public IP when it becomes available.

\end{enumerate}

\newpage

\section{Log into your instance}
Using your key file, ssh into your instance. The username is \texttt{ubuntu}. The process for logging in varies with your operating system and ssh client.  For example, on Linux or OS X, open a terminal window and enter

\begin{verbatim}
  ssh -i <key file name> ubuntu@<ip address>
\end{verbatim}

Once you log into your server, use apt to apply updates and install the \texttt{apache2} package.  Then verify that you can see the default web page by visiting your instance's ip address with a browser.

\section{Wrapping up}
Once you complete the lab, stop your server from the AWS control panel.
\end{document}